% Generated draft: requires mathematical and source review.
\documentclass[11pt]{article}
\usepackage[margin=25mm]{geometry}
\usepackage{fontspec}
\setmainfont{Latin Modern Roman}
\usepackage{amsmath,amssymb,mathtools,unicode-math}
\setmathfont{Latin Modern Math}
\usepackage{enumitem,needspace,xurl,hyperref}
\setlength{\parindent}{0pt}
\setlength{\parskip}{5pt}
\setlength{\emergencystretch}{3em}
\newcommand{\N}{\mathbb{N}}
\newcommand{\Z}{\mathbb{Z}}
\newcommand{\Q}{\mathbb{Q}}
\newcommand{\R}{\mathbb{R}}
\newcommand{\C}{\mathbb{C}}
\newcommand{\F}{\mathbb{F}}
\newcommand{\A}{\mathbb{A}}
\newcommand{\PP}{\mathbb{P}}
\newcommand{\E}{\mathbb{E}}
\newcommand{\eps}{\varepsilon}
\newcommand{\dd}{\,\mathrm{d}}
\begin{document}
\section*{UnsolvedMath ID: 11000002}
\textbf{Source problem code:} catalogue entry
\subsection{Definitions and mathematical statement}
In the usual closed-surface convention, let $S_g$ be a closed connected oriented surface of genus $g$ and define its Torelli group by
\[
 \mathcal I_g=\ker\bigl(\operatorname{Mod}(S_g)
 \longrightarrow\operatorname{Sp}(H_1(S_g;\Z))\bigr).
\]
The excerpt writes $I_g$ without restating the surface convention; boundary or marked-point variants must not be substituted silently. Interpret $H_k(\mathcal I_g)$ as integral group homology unless another coefficient system is specified. A group $G$ is of type $F_k$ if it has an Eilenberg--Mac Lane CW complex $K(G,1)$ with finite $k$-skeleton; $F_1$ is finite generation and $F_2$ finite presentability.

Problem 2.1 asks to determine the finiteness properties of $\mathcal I_g$: for which pairs $(g,k)$ is $H_k(\mathcal I_g;\Z)$ finitely generated as an abelian group, and for which pairs is $\mathcal I_g$ of type $F_k$? These are distinct questions; finite generation as a symplectic-group module is also different from finite generation as an abelian group.

The registry reports the following progress, subject to source verification: $\mathcal I_2$ is infinitely generated free (Mess, 1992); $\mathcal I_g$ is finitely generated for $g\geq3$ (Johnson, 1983), with genus one trivial in the closed convention. For $g\geq2$, it reports $\operatorname{cd}(\mathcal I_g)=3g-5$ (Bestvina--Bux--Margalit, 2007), hence vanishing homology for $k\geq3g-4$, and infinitely generated integral homology for $2g-3\leq k\leq3g-5$ (top degree from Bestvina--Bux--Margalit, extended range attributed to Gaifullin). It also reports finite-dimensional $H_2(\mathcal I_g;\Q)$ for $g\geq51$ and finite generation of $H_2(\mathcal I_g)$ as a symplectic module. Neither statement settles integral finite generation. The summary reports finite presentability for $g\geq3$ still open and separately attributes an infinite-generation statement for $H_*(\mathcal I_g;\Z)$ at $g\geq7$ to Akita, without precisely specifying the latter's algebraic meaning.
\subsection{Short English statement}
Determine, by genus and degree, which integral Torelli homology groups are finitely generated and which finite-skeleton properties the Torelli groups satisfy. Rational homology and symplectic-module finiteness do not settle the integral question.
\subsection{Sources}
\begin{itemize}
\item Local registry: \url{data/unsolvedmath/problems.json}, permanent UnsolvedMath ID $11000002$, exact source code \texttt{catalogue entry}.
\item Farb, \emph{Problems on Mapping Class Groups and Related Topics}, edited volume, Problem 2.1, PDF page 12: \url{https://www.math.uchicago.edu/~farb/papers/mcgbook.pdf}. Supplied access date: 2026-07-29; extraction: pdf-font-structure; 329 labeled Problem, Conjecture and Question headings, disambiguated by page.
\item Supplied progress attributions: Mess (1992), Johnson (1983), Bestvina--Bux--Margalit (2007), Gaifullin and Akita, without complete references. Supplied identifiers: \url{https://arxiv.org/abs/2307.07082} (2023) for rational second homology, and Church--Ershov--Putman, \url{https://arxiv.org/abs/1807.01338}, for symplectic-module finiteness.
\item The summary associates finite presentability with Mess/Birman and Kirby's list, and cross-refers to later chapter Problems 5.12 and 3.1, source codes \texttt{AMR-109-0058} and \texttt{AMR-109-0305}; these attributions and cross-references are retained as supplied.
\end{itemize}
\end{document}
